%%%:%%%%%%%%%%%%%%%%%%%%%%%%%%%%%%%%%%%%%%%%%%%%%%%%%%%%%%%%%%%%%%%%%%%%%%%%%%%%%
%%%: Este é o modelo dun artigo simplificado. NON ESTÁ COMPLETO. ESTO NO É A
%%%: REVISTA FINAL. A revista final componse de varios arquivos que inclúen o que
%%%: aquí é o corpo do documento (dentro do \begin{document}). Esto no meu sistema
%%%: compílase nuns 10 s. A revista_001 completa, en case 50 s
%%%:
%%%: Para usalo fai falta poñer 3 cousas no mesmo directorio
%%%: a clase da revista:                   revista.cls
%%%: este artigo:                          modelo_artigo.tex
%%%: as tipografías:                       fontes\LatinModern
%%%: o esquema da bibliografía             american-physics-society.csl
%%%: a bibliografía                        bibliografia.bib
%%%: (OPCIONAL) configuración de latexmk   latexmkrc
%%%:
%%%: COMO COMPILAR ESTO:
%%%:
%%%: Primeiro, é importante lembrar que os compiladores de TeX/LaTeX son de "unha
%%%: pasada", polo que certa información, como as referencias internas, non se
%%%: xeran ao compilar o documento unha vez. En xeral é necesario darlle varias
%%%: pasadas ao documento, dependendo dos paquetes que usemos. Os comandos exactos
%%%: usados para compilar orixinalmente a revista 002, eran (creo):
%%%:
%%%: lualatex --file-line-error --interaction=nonstopmode --shell-escape  -recorder -output-directory=".aux"  "revistas/002/revista_002.tex"
%%%: biber  ".aux/revista_002.bcf"
%%%: lualatex --file-line-error --interaction=nonstopmode --shell-escape  -recorder -output-directory=".aux"  "revistas/002/revista_002.tex"
%%%: biber  ".aux/revista_002.bcf"
%%%: lualatex --file-line-error --interaction=nonstopmode --shell-escape  -recorder -output-directory=".aux"  "revistas/002/revista_002.tex"
%%%: lualatex --file-line-error --interaction=nonstopmode --shell-escape  -recorder -output-directory=".aux"  "revistas/002/revista_002.tex"
%%%:
%%%: Evidentemente é un rollo escribir esto cada vez. A maneira REAL que usamos
%%%: para compilar a revista é usando o programa 'latexmkrc', o cal executa os
%%%: compiladores e outros executables na orde correcta as veces que faga falta.
%%%: No directorio temos o arquivo 'latexmkrc', que ten unha pequena configuración
%%%: para que 'latexmk' se execute como queremos. Na terminal, simplemente
%%%: escribimos 'latexmk artigo_simplificado.tex' e xa estaría.
%%%:
%%%: Se non queredes traballar a pelo na terminal, sabede na maioría de editores
%%%: debería bastar con premer calquer botón de 'executar'. Se non funciona
%%%: seguramente teñades que cambiar o método que usa voso editor detrás de
%%%: cámaras. Recordo que TeXStudio tiña un panel que deixaba seleccionar que
%%%: compiladores/executables queremos usar. O mellor é usar 'latexmk', o cal vai,
%%%: á súa vez, elixir qué compiladores/executables se usan de xeito automático,
%%%: como xa comentei antes. Nalgúns editores este tipo de compilación pode
%%%: chamarse 'rápida' ou 'automática'.
%%%:
%%%: Todos os paquetes necesarios deberían estar dispoñibles no teu sistema se tes
%%%: unha instalación completa de TeX/LaTeX. En Linux Ubuntu está texlive
%%%: instalable co paquete 'texlive-full' a través de 'apt'. En Windows está
%%%: MikTeX, e en Mac xuraría que existe MacTeX. LaTeX é un rollo de manter en
%%%: distintas plataformas, e en principio só prometemos que funcione en Linux ca
%%%: última versión de texlive. Tamén debería funcionar en Overleaf sen problema
%%%: porque Overleaf usa texlive. Se usades algunha outra plataforma e vos da
%%%: algún erro comentádenolo.
%%%:
%%%: Adicionalmente, fixádevos que este arquivo ten todos os comentarios escritos
%%%: con '###:' Esto é a posta, xa que temos un script aparte para eliminalos de
%%%: golpe cando os recibimos cos contidos dun artigo real. Entón, os VOSOS
%%%: comentarios, escribídeos normal, cun so '#', senón igual o borramos sen querer
%%%:%%%%%%%%%%%%%%%%%%%%%%%%%%%%%%%%%%%%%%%%%%%%%%%%%%%%%%%%%%%%%%%%%%%%%%%%%%%%%
%%%:
%%%: Esto importa a clase definida no arquivo aparte, 'revista.cls'. Usamos a opción
%%%: 'simple' porque non queremos portada, índice nin contraportada neste caso
\documentclass[simple]{revista}
%%%:
\begin{document}
%%%:
%%%: Este macro hai que usalo sempre. Fai bastantes cousas, como engadir o título
%%%: e autor ao índice, reiniciar os conteos de figuras, e evidentemente imprimir
%%%: no PDF a propia información de cada artigo.
%%%:
%%%: \Titular   -> sen asterisco, non engade a sección ao índice
%%%: \Titular*  -> con asterisco, engade a sección ao índice.
%%%:
%%%: Ao ser unha versión simplificada da revista, da igual cal usedes. Eso xa se
%%%: cambiará logo na propia edición
%%%:
%%%: {TEMA}      A escoller entre:
%%%:                divulgacion,historia,actualidadeFacultade,
%%%:                actualidadeCientifica,filosofia,entrevistas,
%%%:                programacion,pasatempos,anuncios,profesorado
%%%: {TITULO}    Evidente, o titulo do artigo. Non debería ser moi longo.
%%%: {AUTORÍA}   Autoría do artigo
%%%: {SUBTITULO} Un pequeno texto algo máis explicativo (e posiblemente longo) co
%%%:             titulo. E o único argumento do titular que pode quedar en branco
\Titular%
{divulgacion}%
{TITULO}%
{AUTORÍA}%
{SUBTITULO}%


%%%: O CORPO DO ARTIGO COMEZA DE AQUÍ EN DIANTE
%%%:%%%%%%%%%%%%%%%%%%%%%%%%%%%%%%%%%%%%%%%%%%%%%%%%%%%%%%%%%%%%%%%%%%%%%%%%%%%%%
%%%:vvvvvvvvvvvvvvvvvvvvvvvvvvvvvvvvvvvvvvvvvvvvvvvvvvvvvvvvvvvvvvvvvvvvvvvvvvvvv
%%%:
%%%:
%%%: DÚAS COLUMNAS ==============================================================
%%%: 'multicols' danos bastante flexibilidade para colocar as cousas en dúas
%%%: columnas. Se unha parte non debe ter 2 columnas, solo hai que rematar o
%%%: entorno 'multicols' antes, escribir o que queiramos (quedará en 1 columna), e
%%%: volver iniciar outro entorno 'multicols' para o resto.
\begin{multicols}{2}

%%%: SECCIÓNS E ESTRUCTURA ======================================================
%%%: Os únicos "niveis" que se permiten son de 'subsection' en abaixo. En
%%%: principio poden estar numerados ou non, a gusto de cada un.
\subsection*{Introdución}

Texto texto texto texto texto texto texto texto texto texto texto texto texto
texto texto texto texto texto texto texto texto

Isto é unha referencia \cite{clave_1919}.
\subsection*{Introdución}

Texto texto texto texto texto texto texto texto texto texto texto texto texto
texto texto texto texto texto texto texto texto

%%%: IMAXES =====================================================================
%%%: Pódense engadir imaxes co entorno 'figure' e ca opción [H] (é obrigatoria)
%%%: Se podedes, intentade que as imaxes non sexan demasiado pesadas, por favor.
\begin{figure}[H] 
    \centering
%%%: Metese a figura como sempre. Pódese xogar co ancho pa axustala.
    \includegraphics[width=0.65\linewidth]{exemplo_imaxe.jpeg}
    \caption{Texto a pe de paxina, interesante...}
    \label{im:exemplo}
\end{figure}

%%%: CITAS TEXTUAIS =============================================================
%%%: Para poñer texto como unha cita 'narrativa', con texto algo máis estreito, en
%%%: itálica, cun nome ao final. Creamos un macro porque todo o mundo as facía
%%%: diferentes, e ás veces funcionaban e ás veces non. Poden usarse como:
\Cita{
A álxebra é xenerosa, frecuentemente da máis do que se lle pide
}
{D'Alembert}

%%%: CITAS E REFERENCIAS BIBLIOGRÁFICAS. ========================================
%%%: O manexo da bibliografía en LaTeX é unha experiencia necesariamente mala,
%%%: sempre. O primeiro punto é asimilalo. Para poder citar cousas, é necesario
%%%: ter un arquivo ao lado (no mesmo directorio) que se chame bibliografía.bib, o
%%%: cal ten a información dos libros, artigos, etc. O contido de dito arquivo é
%%%: algo así (si, un formato tremendamente feo):
%%%:
%%%: @Article{caratheodory_1909,
%%%:   Title    = {Untersuchungen über die Grundlagen der Thermodynamik},
%%%:   Volume   = 67,
%%%:   ISSN     = {0025-5831, 1432-1807},
%%%:   DOI      = {10.1007/BF01450409},
%%%:   Number   = 3,
%%%:   Journal  = {Mathematische Annalen},
%%%:   Author   = {Carathéodory, C.},
%%%:   Year     = 1909,
%%%:   Month    = sep,
%%%:   Pages    = {355–386},
%%%:   Language = {de}
%%%: }
%%%:
%%%: Sen dito arquivo, non se pode xerar a bibliografía automaticamente. Este
%%%: formato de texto dise que está en formato 'bib' ou 'bibtex'. É posible buscar
%%%: ditas referencias en formato 'bib' en internet. En xeral, nas páxinas oficias
%%%: das editoriais que teñen os artigos ou libros, adoita haber unha opción tipo
%%%: 'download citation in bib format' ou algo similar.
%%%:
%%%: Unha vez tendes o arquivo 'bibliografia.bib' nun documento no mesmo
%%%: directorio que esto, pódese citar moi fácil simplemente escribindo (para o
%%%: caso da referencia de Caratheodory de exemplo):
%%%:
%%%: \cite{caratheodory_1909}
%%%:
%%%: Existen outros comandos como \parencite, \textcite, etc. Simplemente mostran
%%%: diferente as citas no propio texto. TENDE COIDADO, PODE QUE NON FUNCIONEN.
%%%: No caso de que queirades que unha referencia apareza ao final de todo, pero non
%%%: queredes citala no propio texto, hai que usar:
%%%:
%%%: \nocite{caratheodory_1909}
%%%:
%%%: Para mostrar a bibliografía ao final, usamos este comando
\printbibliography
%%%:
\end{multicols}
%%%:
\end{document}
